\documentclass[10pt,a4paper]{amsart}
\usepackage[margin=19mm]{geometry}
\usepackage{amsmath,amssymb,mathtools}
\usepackage[scr=rsfs,cal=euler]{mathalfa}
\usepackage{xfrac}
\usepackage[unicode,hidelinks,bookmarks=false]{hyperref}
\newcommand{\ie}{i.e.\ }
\newcommand{\cf}{cf.\ }
\newcommand{\resp}{respectively\ }
\newcommand{\Subsection}{\subsection}
\providecommand{\nobreakdash}{\nobreak\hbox{-}}
\setlength{\emergencystretch}{2em}
\multlinegap0pt
\usepackage{bm}
\usepackage{bbm}
\usepackage{dsfont}
\usepackage{xspace}
\usepackage{cancel}
\usepackage{mathtools}
\usepackage{amsthm}
\newtheorem{lemma}{Lemma}
\newcommand{\sumstar}{\sideset{}{^*}\sum}
\title[Second moment of $\textrm{GL(3)} \times \textrm{GL(2)}$ $L$-]{Second moment of $\textrm{GL(3)} \times \textrm{GL(2)}$ $L$--functions}
\author{Sumit Kumar, K. Mallesham, Suraj Panigrahy}
\date{Source: arXiv:2602.22127v1 (CC BY 4.0)}
\begin{document}
\maketitle

\noindent\emph{Source and licence.} Second moment of $\textrm{GL(3)} \times \textrm{GL(2)}$ $L$--functions. Version arXiv:2602.22127v1 (CC BY 4.0), distributed under the Creative Commons Attribution 4.0 International licence, as recorded in the arXiv licence metadata for this exact version and shown on the abstract page https://arxiv.org/abs/2602.22127v1. Full rights notice: licenses/arxiv-2602.22127-CC-BY-4.0.txt.\par\smallskip

\noindent\textit{Editorial scope.} Counted unit: the Lemma whose source label is \texttt{hybrid large-sieve with level} in the article (source0.tex lines 400--405, source label \texttt{hybrid large-sieve with level}), delivered together with the article's own complete original proof of it (source0.tex lines 406--480, one \texttt{proof} environment of 75 lines). No mathematics is written, repaired or re-derived: statement and proof are the article's own text line for line. Required context retained uncounted: large sieve with level (lines 276--289). Every definition, display and label of those lines is retained unchanged. Omitted from this package and not redistributed: 1--275, 290--399, 481--1601. Nothing inside a retained excerpt is omitted or altered: every retained body is delivered exactly as the article's lines are written, so no omission receipt of any kind accompanies this package. Citations: the retained excerpts carry no citation command at all, so no bibliography entry of the article is transcribed and no reference is redistributed. Reference adaptation: none was needed and none was made; every reference inside the delivered unit resolves inside it, so no unresolved reference or citation is produced. Printed slips kept literal: no printed word, symbol, hypothesis or number of the counted statement or of its proof was altered. Layout and class: the article's own class is \texttt{amsart} with packages including amsfonts, amsthm, amsmath, amssymb, mathrsfs, hyperref, blindtext, geometry, helvet; the wrapper is the lane's pinned \texttt{amsart} editorial wrapper at 10pt on A4, which restates the theorem-like environments the retained text needs and adds the article's own macro definitions that the retained text needs (\texttt{\textbackslash sumstar}), with extra wrapper packages (bm, bbm, dsfont, xspace, cancel, mathtools). The delivered numbering is the unit's own. Assets: no retained span contains a figure environment, a graphics-inclusion command, a PDF-page inclusion, a table or a TikZ picture; no image file is copied into this package, the package declares no asset dependency and no image file is redistributed with it. Licence: version arXiv:2602.22127v1 of this article is distributed under the Creative Commons Attribution 4.0 International licence, as recorded in the arXiv licence metadata for this exact version and shown on the abstract page https://arxiv.org/abs/2602.22127v1. A full rights notice is delivered at licenses/arxiv-2602.22127-CC-BY-4.0.txt. Characters: every retained excerpt is pure ASCII, so the delivered engine, XeLaTeX with its default Latin Modern text font, reports no missing character.
\begin{lemma}\label{large sieve with level}
    Let $a_m$ be any complex numbers and $D>0$ be a fixed  integer. Let $1 \leq M\leq N$. Then  we have 
    \begin{align}\label{small N}
   S(C,N,M;D):= \mathop{\sum_{\substack{c \sim C \\ (c,D)=1}}}\sumstar_{\gamma (cD)}  \Big|\sum_{N \leq m < N+M} a_m  e\left( \frac{\gamma m}{cD} \right) \Big| ^2  \ll  C^\epsilon(M +  DC^2) \sum_m|a_m|^2.
\end{align}
%where the second factor occurs only if $DC^2D^\epsilon \geq NQ$. In particular for $Q=D^{5 \epsilon}$ and $DC^2D^\epsilon < NQ$, we have 
%\begin{align}\label{only zero freq}
   %  S(C,N;D) \ll N D^{5 \epsilon} \sum_m|a_m|^2. 
%\end{align}
%We also have 
%\begin{align}
  %  S(C,N;D) \leq \mathop{\sum_{\substack{c \sim C \\ (c,D)=1}}}\sum_{\gamma (cD)}  \Big|\sum_{N \leq m <  2N} a_m  e\left( \frac{\gamma m}{cD} \right) \Big| ^2 \leq (5C^2D+NC) \sum_m |a_{m}|^2.
%\end{align}
\end{lemma}

\begin{lemma}\label{hybrid large-sieve with level}
    Let $a_m$ be  any complex numbers  and $D>0$ be a fixed  integer and $T >0 $ be a real number. Let $f$ be a $C^1$ function on $[N,2N]$ such that $f^\prime$ does not vanish. Let  $X=\sup_{x \in [N,2N]} \frac{1}{|f^\prime(x)|}$. Then we have  
    \begin{align*}
   S_T(C,N;D):=\mathop{\sum_{\substack{c \sim C \\ (c,D)=1}}}\sumstar_{\gamma (cD)} \int_{-T}^{T} \Big|\sum_{N \leq m <  2N} a_m  e(tf(m)) e\left( \frac{\gamma m}{cD} \right) \Big| ^2\mathrm{d}t  \ll  (C^2DT+X) \sum_m|a_m|^2.
\end{align*}
\end{lemma}

\begin{proof}
Let $W$ be a non-negative smooth bump  function supported on $[1/2,3/2]$ such that $W=1$ on $[1,2]$. Thus we have 
  \begin{align*}
   S_T(C,N;D) \leq  \mathop{\sum_{\substack{c \sim C \\ (c,D)=1}}}\sumstar_{\gamma (cD)} \int_{0}^{\infty}W(t/T) \Big|\sum_{N \leq m <  2N} a_m  e(tf(m)) e\left( \frac{\gamma m}{cD} \right) \Big| ^2\mathrm{d}t.
\end{align*}
For any sequence of complex numbers $(b_m)$, we have 
\begin{align}\label{open the abs}
    \int_{T}^{2T } \Big|\sum_{N \leq m <  2N} b_m  e(tf(m))  \Big| ^2\mathrm{d}t \leq  \sum_{m_1}  \sum_{m_2}b_{m_1} \overline{b}_{m_2} \int_{0}^{\infty} W\left( \frac{t}{T} \right) e\left( t(f(m_1)-f(m_2) \right)  \mathrm{d}t 
\end{align}
 %Opening  the absolute valued square  we arrive at 
 %\begin{align*}
 %    \mathop{\sum_{\substack{c \sim C \\ (c,D)=1}}} \sumstar_{\gamma (cD)}   \sum_{m_1}  \sum_{m_2}a_{m_1} \overline{a}_{m_2} e\left(\frac{ (m_1-m_2)\gamma  }{Dc} \right) \int_{0}^{\infty} W(\frac{t}{T}) e\left( t(f(m_1)-f(m_2) \right)  \mathrm{d}t 
 %\end{align*}
 By a change of variable, we see that the $t$-integral is negligibly small unless $$ |f(m_1)-f(m_2)| \leq C^{\epsilon}/T. $$ 
 By the mean-value theorem, we have
 \[|m_2-m_1| \leq {|f(m_1)-f(m_2)|X} \leq {C^\epsilon X}/{T}.\]
Let $Y=\min\{C^\epsilon X/T,N \}$. We dissect the sum over $m_1$ and $m_2$ into the intervals $$I_k=[N+kY,N+(k+1)Y), \ 0 \leq k \ll N/Y $$   as follows:
\[    \sum_{N \leq m_1<2N}  \sum_{N \leq m_2 <2N}\cdots = \sum_{0\leq k\ll N/Y}\sum_{m_1 \in I_k} \sum_{m_2 \in J_k} \cdots,\]
where $J_k=I_k+([-Y,Y]\cap \mathbb{Z})\cap [N,2N).$
Thus the right side of \eqref{open the abs} is 
\begin{align*}
   &\leq \frac{1}{2} \int_{0}^{\infty} W\left( \frac{t}{T} \right) \sum_k \Big| \sum_{m_1 \in I_k} b_{m_1} e\left( t(f(m_1) \right) \Big|^2+\Big|\sum_{m_2 \in J_k} \bar{b}_{m_2} e\left( -tf(m_2) \right) \Big|^2\mathrm{d}t.
     \end{align*}
We focus on the second term, as the first term is similar, which, on taking  $b_m=a_m e\left(\frac{\gamma m}{cD}\right)$ is given by 
\begin{align*}
    \frac{1}{2} \int_{0}^{\infty} W\left( \frac{t}{T} \right) \sum_k \Big|\sum_{m \in J_k} a_m e\left(\frac{\gamma m}{cD}\right) e\left( tf(m) \right) \Big|^2\mathrm{d}t.
\end{align*}
Thus on summing over $c$ and $\gamma$, we arrive at 
\begin{align*}
     \frac{1}{2} \int_{0}^{\infty} \mathop{\sum_{\substack{c \sim C \\ (c,D)=1}}}\sumstar_{\gamma (cD)} W\left( \frac{t}{T} \right) \sum_k \Big|\sum_{m \in J_k} a_m e\left(\frac{\gamma m}{cD}\right) e\left( tf(m) \right) \Big|^2\mathrm{d}t.
\end{align*}
On applying Lemma \ref{large sieve with level} with the coefficients $a_m e(t(f(m)))$, the above is  
\[ \ll C^\epsilon \sum_{k \ll N/Y} T (C^2D+ Y)\sum_{m \in J_k}|a_m|^2 \ll C^\epsilon (TC^2D + X )\sum_m|a_m|^2.\]
The lemma follows.
 % Like  Lemma \ref{large sieve with level}, this lemma is also a consequence of duality principle and Poisson summation.   We apply the duality lemma (see Lemma \ref{dualitylemma}) to $ S_T(C,N;D) $  to  see that  
%\begin{align*}
 %    S_T(C,N;D)
  %     \ll  \sum_m|a_m|^2 \sup_{\|\beta\|_2=1} S_T^\beta(C,N;D) ,
%\end{align*}
%where 
%\begin{align}
   %   S_T^\beta(C,N;D) =\sum_{N \leq m \leq 2N} \Big|    \mathop{\sum_{\substack{c \sim C \\ (c,D)=1}}}\sumstar_{\gamma (cD)} \int_{-T}^{T} \beta(\gamma,t,c)  e\left(\frac{ m\gamma }{cD} \right)  e\left(f(m,c) \right) m^{it}\mathrm{d}t \Big|^2.
%\end{align}
 %Next 
 %smooth out  the sum over $m$ with a bump function $w(x)$ such that $0\leq w(x) \leq 1$,  $w(x)=1$ for $x \in [1,2] $ and $w(x)=0$, for $x \notin (1/2,3/2)$,  $x^j w^{(j)}(x) \ll_{j}1$.  Opening the absolute valued square, we arrive at
%\begin{align*}
   %  S_T^\beta(C,N;D) =&    \mathop{\sum_{\substack{c \sim C \\ (c,D)=1}}} \mathop{\sum_{\substack{c^\prime \sim C \\ (c^\prime,D)=1}}} \sumstar_{\gamma (cD)} \sumstar_{\gamma (c^\prime D)} \int_{-T}^{T} \int_{-T}^{T}\beta(\cdots) \overline{\beta}(\cdots)  \\
   %  & \times \sum_m w(m/N) e\left(\frac{ m(\gamma c^\prime-\gamma^\prime c ) }{Dcc^\prime} \right)  e\left( f(m,c)-f(m,c^\prime) \right) m^{ {i(t-t^\prime)}{}} \mathrm{d}t \,\mathrm{d}t^\prime.
%\end{align*}
%On applying the Poisson summation formula with  modulus $Dcc^\prime$, we see that  the sum over $m$ is given by
%\begin{align}
 %   \sum_m \cdots= {N}{ }\sum_{m  }   \delta( \gamma c^\prime-\gamma^\prime c \equiv -m\, \mathrm{mod}\, Dcc^\prime)  I(m,c, c^\prime,t,t^\prime),
%\end{align}
%where 
%\begin{align}
   % I(\cdots)= \int w(z)   e\left( f(Nz,c)-f(Nz,c^\prime) \right) (Nz)^{{i(t-t^\prime)}{}} e\left( \frac{-mNz}{Dcc^\prime}\right) \mathrm{d}z
%\end{align}
%An integration by parts argument shows that $ I(\cdots)$ is negligibly small unless 
%$$ |m| \leq N_0:=D^{\epsilon}\left( 1+ T
%\right) \frac{DC^2}{N}. $$
%For $N \geq D^{2 \epsilon} TDC^2$, we note that  $N_0<1$ and  thus  $m=0$. We analyse this now. 
% For $m=0$, the congruence condition yields
% $c=c^\prime$,  $ \gamma= \gamma^\prime (c),$ and 
% $\gamma= \gamma^\prime (D)$. Moreover, the integral transform  $I(\cdots)$ ensures $|t-t^\prime| \leq D^{\epsilon}$ at the cost of a negligible error term. Indeed, in the range $|t-t^\prime| \gg D^{\epsilon}$,  $I(\cdots) = O(D^{-A})$ using integration by parts).  
%In the range $|t-t^\prime| \leq D^{\epsilon}$ we use the trivial bound    $ I(\cdots) \ll 1 $. Thus   we get 
% \begin{align*}
 %    S_T^\beta(C,N;D)  \ll N &   \mathop{\sum_{\substack{c \sim C \\ (c,D)=1}}}  \ \sumstar_{\substack{ \gamma \shortmod{cD}}  }  \int_{-T}^T \int_{-T}^T  |\beta(\gamma,t,c) \overline{\beta}( \gamma, t^\prime, c)| \delta(|t-t^\prime| \leq D^\epsilon)  \mathrm{d}t \, \mathrm{d}{t^\prime}.
%\end{align*}
%By a  change of variable $t-t^\prime=t_1$, $|t_1| \leq D^\epsilon$, we see that 
%\begin{align*}
%     S_T^\beta(C,N;D)  &\ll N     \mathop{\sum_{\substack{c \sim C \\ (c,D)=1}}}  \ \sumstar_{\substack{ \gamma \shortmod{cD}}  }  \int_{-T}^T \int_{-T}^T  |\beta(\gamma,y_1+y^\prime,d) \overline{\beta}( \gamma, y^\prime, d)| \mathrm{d}y_1 \, \mathrm{d}{y^\prime} \\
  %   &\ll N     \mathop{\sum_{\substack{c \sim C \\ (c,D)=1}}}  \ \sumstar_{\substack{ \gamma \shortmod{cD}}  }  \int_{-T}^T \int_{-T}^T (|\beta(\gamma,y_1+y^\prime,d)|^2+ |\overline{\beta}( \gamma, y^\prime, d)|^2) \mathrm{d}y_1 \, \mathrm{d}{y^\prime} \ll N.
%\end{align*}
%Thus the lemma follows.
\end{proof}

\end{document}
